\clearpage\thispagestyle{empty}
\pdfbookmark[0]{Quellen zur Besprechung von Müllers Anleitung zu Sonnenfinsternissen}{mueller1816-sources}
\begin{center}\Large\bfseries Quellen zur Besprechung von Müllers Anleitung zu Sonnenfinsternissen\end{center}

Die Übertragung folgt Carl Friedrich Gauß’ Besprechung in den \emph{Werken}, Bd.~VI (Göttingen, 1874), S.~591. Der Erstabdruck erschien in den \emph{Göttingischen gelehrten Anzeigen}, Stück~20, 3.~Februar 1816, S.~200. Die Form \emph{merkwürdiegen} steht so in der Werkausgabe.

Johann Wolfgang Müllers Anleitung erschien 1816 in Sulzbach bei J.~E. Seidel. Das digitalisierte Exemplar, Astr.p.~94~r, stammt aus der Bayerischen Staatsbibliothek. In der abschließenden Gegenüberstellung auf S.~128 gibt Müller die Größe der Finsternis in den historischen Einheiten Zoll und Minuten an.

KI-Übertragung aus dem historischen Deutsch in modernes Deutsch. KI-erstellter Satz.

\href{https://gdz.sub.uni-goettingen.de/id/PPN236007467}{Werkausgabe}. \href{https://images.p.sub.uni-goettingen.de/iiif/image/gdz:PPN319721507_1816:00000202/full/1650,/0/default.jpg}{Erstdruck}. \href{https://mdz-nbn-resolving.de/urn:nbn:de:bvb:12-bsb10060601-5}{Müllers Buch}. \href{https://api.digitale-sammlungen.de/iiif/image/v2/bsb10060601_00130/full/1500,/0/default.jpg}{Gegenüberstellung auf S.~128}. \href{https://api.digitale-sammlungen.de/iiif/image/v2/bsb10060601_00131/full/1500,/0/default.jpg}{Kupfertafel}.

\clearpage
